Solving the characteristic equation $(r-2)(r+1)=0$, or checking by induction, gives
\[
a_n=\dfrac{2^n-(-1)^n}{3}.
\]
Since $3$ is invertible modulo $61$, the divisibility is equivalent to $(-2)^k\equiv1\pmod{61}$. Fermat's theorem shows its order divides $60$. Direct modular powers give
\[
(-2)^{30}\equiv-1,\qquad (-2)^{20}\equiv47,\qquad (-2)^{12}\equiv9\pmod{61}.
\]
None is $1$. Every proper divisor of $60$ divides at least one of $30,20,12$, so the order is $60$. Hence $\boxed{k=60}$.
